% Generated from locale/tr/content/history/biographies/julia-robinson.tex; do not hand-edit.


\olfileid[tr]{his}{bio}{rob}

\olsection{Julia Robinson}

\olphoto{robinson-julia}{Julia Robinson}

Julia Bowman Robinson Amerikalı bir matematikçiydi. Başlıca karar problemleri
üzerine çalışmalarıyla, en çok da Hilbert'in onuncu probleminin çözümüne
katkılarıyla tanınır. Robinson 8 Aralık 1919'da Missouri, St.~Louis'de doğdu.
Daha çocukken sayılara ilgi duyduğunu anlatır \citep[4]{Reid1986}. Dokuz
yaşında kızıl hastalığına yakalandı ve birkaç kez tekrarlayan romatizmal ateş
nöbetleri geçirdi. Bu nedenle zamanının büyük bölümünü yatakta geçirmek
zorunda kaldı ve eğitiminde geri düştü. Özel öğretmenlerin yardımıyla arayı
kapatabilse de hastalığının fiziksel etkileri yaşamı boyunca sürdü.

Çocukluğunda yaşadığı güçlüklere karşın Robinson liseyi matematik ve fen
alanında birçok ödülle bitirdi. Üniversite öğrenimine San Diego State
College'da başladı; son sınıf öğrencisiyken Berkeley'deki California
Üniversitesi'ne geçti. Burada matematikçi Raphael Robinson'dan etkilendi.
İkisi yakın arkadaş oldu ve 1941'de evlendi. Bir öğretim üyesinin eşi olduğu
için Robinson'ın Berkeley matematik bölümünde ders vermesi yasaktı. Matematik
derslerini dinleyici olarak izlemeyi sürdürse de üniversiteden ayrılıp bir
aile kurmayı umuyordu. Ne var ki düğününden kısa süre sonra zatürreye
yakalandı. Kendisine, çocukken geçirdiği romatizmal ateş nedeniyle kalbinde
önemli miktarda nedbe dokusu biriktiği söylendi. Doktor, dokunun ağırlığı
yüzünden kırk yaşını geçemeyeceğini öngördü ve çocuk sahibi olmamasını
öğütledi \citep[13]{Reid1986}.

Robinson uzun süre depresyondaydı; fakat sonunda matematik öğrenimini
sürdürmeye karar verdi. Berkeley'e döndü ve Alfred Tarski'nin danışmanlığında
doktorasını 1948'de tamamladı. Gerçek sayıların birinci dereceden teorisinin
karar verilebilir olduğu Tarski tarafından gösterilmişti; G\"odel'in
çalışmasından da doğal sayıların birinci dereceden teorisinin karar
verilemezliği çıkıyordu. Rasyonel sayıların birinci dereceden teorisinin
karar verilebilir olup olmadığı önemli bir açık problemdi. Robinson tezinde
\citeyearpar{Robinson1949} bunun karar verilemez olduğunu ispatladı.

Karar problemleriyle ilgilenen Robinson daha sonra Hilbert'in onuncu
problemine bir çözüm bulmaya girişti. Bu problem, David Hilbert'in 1900'de
ortaya attığı 23 ünlü matematik probleminin biriydi. Onuncu problem, $3x^2 -
2y +3 = 0$ gibi tamsayı katsayılı bir polinom denkleminin tamsayılarda çözümü
bulunup bulunmadığını sonlu bir sürede yanıtlayacak bir algoritmanın var olup
olmadığını sorar. Böyle sorulara \emph{Diofantos problemleri} denir. İlk bazı
başarılarının ardından Robinson, problem üzerinde çalışan Martin Davis ve
Hilary Putnam'la güçlerini birleştirdi. Üstel Diofantos problemlerinin
(bilinmeyenlerin üs olarak da yer alabildiği problemler) karar verilemez
olduğunu göstermeyi başardılar ve belirli bir varsayımın (sonradan ``J.R.''
diye adlandırıldı) Hilbert'in onuncu probleminin karar verilemezliğini
gerektirdiğini gösterdiler \citep{DavisPutnamRobinson1961}. Robinson 1960'lar
boyunca problem üzerinde çalışmayı sürdürdü. Genç Rus matematikçi Yuri
Matijasevich, 1970'te J.R. varsayımını nihayet ispatladı. Birleşik sonuç bugün
Matijasevich--Robinson--Davis--Putnam teoremi, kısaca MRDP teoremi diye
adlandırılır. Matijasevich ile Robinson arkadaş oldular ve birkaç makalede
birlikte çalıştılar. Robinson bir keresinde Matijasevich'e yazdığı mektupta
``aslında (binlerce kilometre uzaktan) birlikte çalışırken ikimizin de tek
başına kaydedebileceğinden daha fazla ilerlediğimiz açıkça görüldüğü için çok
memnunum'' demişti \citep[45]{Matijasevich1992}.

Robinson, Amerikan Matematik Derneği'nin ilk kadın başkanı ve Ulusal Bilimler
Akademisi'ne seçilen ilk kadındı. Lösemi tanısı konduktan sonra 30 Temmuz
1985'te 65 yaşında öldü.

\begin{reading}
Robinson'ın matematik makaleleri \textit{Collected Works} içinde mevcuttur
\citep{Robinson1996}; bu derleme, Ulusal Bilimler Akademisi için yazılmış
biyografik anı yazısının yeniden basımını da içerir \citep{Feferman1994}.
Robinson'ın
ablası Constance Reid, söyleşilere dayanarak ``Autobiography of Julia''yı
\citep{Reid1986} ve ayrıca tam bir anı kitabını \citep{Reid1996} yayımladı.
Robinson ve Hilbert'in onuncu problemi üzerine kısa bir belgeseli George
Csicsery yönetti \citep{Csicsery2016}. Yuri Matijasevich'in Robinson'la iş
birlikleri ve Robinson'ın çalışmaları üzerindeki etkisi hakkında kısa bir anı
yazısı için \citep{Matijasevich1992} kaynağına bakın.
\end{reading}

